%% ═══════════════════════════════════════════════════════════════════════════════
\section{Introduction}
\label{sec:intro}
%% ═══════════════════════════════════════════════════════════════════════════════

Deep neural networks are successful in various tasks but are also susceptible to adversarial examples: malicious input perturbations designed to deceive the network~\cite{szegedy2014,carlini2017,madry2018}. Many adversarial attacks target image classifiers and compute either an imperceptible change, bounded by a small $\varepsilon$ with respect to an $L_p$ norm, or a perceivable change obtained by perturbing a semantic feature of the input, such as brightness, translation, or occlusion.

A popular safety property for understanding a network's resilience to such attacks is \emph{local robustness}. Local robustness is parameterized by an input and its neighborhood: for a network classifier, the goal is to prove that the network classifies all inputs in the neighborhood the same. Several local robustness verifiers have been introduced for checking robustness in an $\varepsilon$-ball~\cite{singh2019,tjeng2019,wang2021}. However, local robustness is limited to reasoning about a single neighborhood at a time. To reason about a network's robustness over \emph{all} possible inputs, VHAGaR~\cite{vaghar} introduced the concept of \emph{global robustness}: computing the \emph{minimal globally robust bound}, the smallest confidence above which the network is provably robust to a given perturbation.

In practice, neural networks are rarely deployed in isolation. A common scenario involves training a more accurate network $N_2$ from an existing network $N_1$ via fine-tuning, distillation, or retraining. A natural question arises: \emph{does the robustness guarantee of $N_1$ transfer to $N_2$?} More precisely, if $N_1$ is known to be globally robust above some confidence threshold $\dN{1}$, can we certify $N_2$'s robustness without re-running the expensive global verification from scratch?

In this work, we address the problem of \emph{robustness transfer verification}. Given a source network $N_1$ and a target network $N_2$, we compute a confidence margin $\dD$ such that any input where $N_2$'s confidence exceeds $N_1$'s by more than $\dD$ inherits $N_1$'s robustness guarantee. Intuitively, $\dD$ captures the worst-case confidence gap between the two networks in the regime where $N_1$ is certifiably robust but $N_2$ can still be fooled. This problem is challenging for two main reasons. First, it requires encoding three forward passes --- $N_1(x)$, $N_2(x)$, and $N_2(\xp)$ --- simultaneously in a single optimization problem. Second, the perturbation introduces complex dependencies between the clean and perturbed computations that must be captured precisely for the resulting bounds to be tight.

We extend VHAGaR with a \emph{transfer mode} that encodes the robustness transfer problem as a mixed-integer program (MIP). Our MIP simultaneously models three forward passes --- $N_1(x)$, $N_2(x)$, and $N_2(\xp)$ --- linked by the perturbation constraint $\xp \in \mI_\ep(x)$. To make the resulting MIP tractable, we introduce four complementary tightening techniques: (1)~a \emph{hyper-adversarial attack} that provides a warm-start solution via projected gradient descent, (2)~\emph{perturbation-difference interval propagation} that bounds the activation differences between clean and perturbed passes, (3)~\emph{inter-network interval bounds} exploiting the structural similarity between $N_1$ and $N_2$, and (4)~\emph{dependency analysis} that identifies monotone relationships between activations. Our approach currently supports the $L_\infty$ perturbation and six semantic feature perturbations.

We provide a complete formal analysis of the transfer MIP's guarantees. We prove that the computed margin $\dD$ provides a \emph{sound lower bound}: $\dN{2} \geq \dN{1} + \dD$, meaning $N_2$'s global robustness bound is at least the sum of $N_1$'s bound and the transfer margin. We then exhibit a concrete counterexample with two simple linear networks demonstrating that the reverse inequality is \emph{false} in general --- the bound is strict. Finally, we propose a modified transfer MIP, replacing the constraint that $N_1$ is robust with the constraint that $N_1$ is \emph{foolable}, and prove that the resulting quantity $\ddiffmod$ satisfies the complementary upper bound $\dN{1} + \ddiffmod \geq \dNmod{2}$.

In summary, our main contributions are:
(1)~formalizing the robustness transfer problem and encoding it as a MIP with four tightening techniques (Section~\ref{sec:transfer}),
(2)~proving that the transfer MIP yields a sound lower bound on $N_2$'s robustness improvement over $N_1$ (Section~\ref{sec:formal}),
(3)~exhibiting an explicit counterexample showing the bound is strict, with a precise characterization of the gap (Section~\ref{sec:formal}), and
(4)~designing a modified formulation that provides a matching upper bound over the joint-foolable domain (Section~\ref{sec:formal}).

The rest of this paper is organized as follows. Section~\ref{sec:notation} provides background on network classifiers and formalizes the central quantities. Section~\ref{sec:perturbations} defines the supported perturbation models. Sections~\ref{sec:standard} and~\ref{sec:transfer} present the standard-mode and transfer-mode MIP formulations, respectively. Section~\ref{sec:modes} discusses the two attack constraint variants. Section~\ref{sec:formal} provides the full formal analysis. Section~\ref{sec:flags} describes the four tightening techniques. Section~\ref{sec:mip} gives the complete MIP formulation, and Sections~\ref{sec:inference} and~\ref{sec:conclusion} present the runtime inference procedure and conclusions.
